
% ---------------------------------------------------------------------------
% Résumé en français et mots-clés
% ---------------------------------------------------------------------------

\chapter*{Résumé}

Le produit Scan d'ISOGEO permet d'explorer différentes sources de données
géographiques afin d'identifier leur contenu et d'en extraire les métadonnées.
Ces informations alimentent ensuite le processus de catalogage, qui constitue
le cœur de l'activité d'ISOGEO en facilitant le recensement, la documentation et
la valorisation des données géographiques. Scan prend notamment en charge des
bases de données, des géodatabases d'entreprise ainsi que plusieurs formats de
fichiers géographiques.

Historiquement, l'ensemble de ces traitements reposait sur des workbenches FME,
ce qui entraînait une dépendance externe et des coûts de licence. ISOGEO a donc
engagé leur remplacement progressif par des scripts Python intégrés directement
au produit. Ma mission a consisté à poursuivre cette migration, déjà initiée pour
certaines sources comme PostGIS et les géodatabases d'entreprise Esri.

Après une phase d'analyse et de documentation de l'existant, j'ai participé au
développement de traitements pour plusieurs sources, notamment les données
vectorielles et raster, GeoPackage, les données CAO, Oracle, SQL Server et
GeoParquet. Les travaux les plus récents ont également porté sur les nuages de
points aux formats LAS et LAZ. Ces scripts permettent d'extraire les tables, les
champs, les projections, les géométries, les emprises spatiales et les autres
métadonnées attendues par Scan, ainsi que de générer une signature des données.

La bibliothèque GDAL/OGR a été principalement utilisée, complétée par des
bibliothèques adaptées à certains formats spécifiques, comme LibreDWG ou encore SQLAlchemy. Les scripts ont ensuite été organisés dans une architecture modulaire et
préparés pour être distribués sous la forme d'un exécutable Windows autonome avec
PyInstaller. Cet exécutable est généré dans un environnement Docker afin de
rendre le processus reproductible et indépendant des bibliothèques installées sur
le poste du développeur. Des tests de régression et des tests d'intégrité de la
signature permettent également de valider les traitements.

Ces travaux contribuent à réduire la dépendance de Scan à FME, à faciliter la
maintenance des traitements et à renforcer leur intégration dans
l'environnement logiciel d'ISOGEO.

\vspace{0.7cm}

\noindent
\textbf{Mots-clés :}
géomatique ; Scan ; ISOGEO ; Python ; migration FME ; métadonnées ; signature ;
GDAL/OGR ; GeoPackage ; GeoParquet ; CAO ; bases de données spatiales ;
nuages de points ; Docker ; PyInstaller.
